
\chapter{Formulación del Proyecto}

\section{Introducción}

El presente documento se ha elaborado en respuesta a la solicitud de [Nombre del Cliente o Empresa Cliente] para el desarrollo de un sistema de base de datos web destinado a gestionar la información académica y de investigación de su institución. Comprendemos la necesidad de contar con una plataforma robusta, escalable y eficiente que permita administrar de manera integrada los datos relativos a profesores, alumnos, proyectos, departamentos y asignaturas.

Este informe tiene como objetivo principal presentar una propuesta detallada del modelo de datos que servirá como pilar fundamental para la construcción del sistema. Un modelo de datos bien definido es crucial para garantizar la integridad de la información, facilitar el desarrollo de funcionalidades complejas, optimizar las consultas y asegurar la futura adaptabilidad del sistema a nuevas necesidades.

La información contenida en las siguientes secciones ha sido cuidadosamente analizada a partir de los requisitos proporcionados. Se expondrá la identificación de las entidades principales, sus atributos y las relaciones que existen entre ellas. Esta estructura ha sido pensada para ser comprensible tanto por el equipo de gerencia, para la toma de decisiones estratégicas, como por el equipo de desarrollo, quienes se encargarán de la implementación técnica.

Confiamos en que este análisis y la propuesta de diseño que se presenta a continuación constituirán una base sólida para el exitoso desarrollo del sistema de base de datos web que su institución requiere. Estamos a su disposición para discutir cualquier aspecto de este informe y para colaborar estrechamente en las siguientes fases del proyecto.

Perfecto. A continuación, te presento una propuesta para el objetivo general y los objetivos específicos del proyecto, considerando las etapas que mencionaste:

\section{Objetivos del Proyecto}

La arquitectura metodológica de este trabajo asume los lineamientos de la iniciativa CDIO, marco que postula un ciclo de vida estructurado para el desarrollo de sistemas ingenieriles[cite: 1]. Considerando los límites operativos del encargo, el proyecto transita exclusivamente por las fases de Concepción, Diseño e Implementación (CDI)[cite: 1], delegando la etapa de Operación comercial y sustituyéndola por una validación funcional a nivel de prototipo computacional.

\subsection{Objetivo General}
Desarrollar un sistema de base de datos web para la gestión de información académica e investigativa en el contexto de [Nombre del Cliente o Empresa Cliente], con el propósito de optimizar el acceso concurrente y centralizar la administración de los registros institucionales.

\subsection{Objetivos Específicos}
\begin{enumerate}
	\item \textbf{Concebir:} Formular los requerimientos técnicos y operacionales mediante el diagnóstico de las necesidades organizacionales del cliente, delimitando la factibilidad del sistema de información.
	\item \textbf{Diseñar:} Modelar la arquitectura de software y los esquemas lógicos de la base de datos relacional que estructurarán la solución tecnológica subyacente.
	\item \textbf{Implementar:} Construir el artefacto informático a nivel de prototipo, integrando los lenguajes de programación web/móvil y los motores de datos seleccionados en un entorno de desarrollo controlado.
	\item Evaluar la consistencia y rendimiento del prototipo funcional mediante simulaciones de carga y pruebas estructurales, verificando las métricas de calidad sin contemplar su despliegue en un entorno operativo de producción.
\end{enumerate}

\section{Metodología de Trabajo}\label{sec:mettrabajo}

El desglose de actividades se estructura en subordinación a los propósitos planteados, estableciendo un flujo lógico desde la abstracción teórica hasta la verificación empírica del código fuente. Esta configuración asegura una trazabilidad coherente a lo largo del ciclo de vida del software, alineando el desarrollo del prototipo con los requerimientos académicos.

\subsection{Objetivo Específico 1: Fase de Concepción y Formulación}
\subsubsection{Actividades y Tareas}
\begin{enumerate}
	\item Diagnóstico de Requerimientos y Extracción de Entidades:
	\begin{enumerate}
		\item Levantar los requerimientos funcionales extrayendo las entidades principales del dominio del problema (Profesor, Proyecto, Alumno, ProgramaPostgrado, Departamento, Asignatura).
		\item Evaluar la factibilidad tecnológica preliminar para la plataforma web, delimitando el alcance del prototipo.
	\end{enumerate}
	\item Consolidación Documental y Planificación Temporal:
	\begin{enumerate}
		\item Estructurar el informe técnico inicial, argumentando la justificación del desarrollo y definiendo la matriz de planificación temporal (Carta Gantt) para las iteraciones.
		\item Revisar y validar la comprensión de los requerimientos, asegurando que el modelo mental concuerda con las necesidades institucionales.
	\end{enumerate}
\end{enumerate}

\subsection{Objetivo Específico 2: Fase de Diseño y Modelado}
\subsubsection{Actividades y Tareas}
\begin{enumerate}
	\item Modelado Conceptual y Lógico de Datos:
	\begin{enumerate}
		\item Proyectar el Modelo Entidad-Relación (MER), definiendo atributos clave (por ejemplo, el RUN como identificador único) y estableciendo las cardinalidades y restricciones (1:N, M:N).
		\item Transformar el MER al modelo relacional lógico, ejecutando un proceso de normalización riguroso hasta la Tercera Forma Normal (3FN) para evitar anomalías transaccionales y dependencias transitivas.
	\end{enumerate}
	\item Arquitectura de Interfaces y Experiencia de Usuario:
	\begin{enumerate}
		\item Elaborar wireframes de bajo nivel para esquematizar la estructura y navegación del sistema.
		\item Construir mockups de alta fidelidad que proyecten la apariencia visual, considerando el diseño diferenciado de interfaces para usuarios con privilegios limitados frente a los administradores.
	\end{enumerate}
\end{enumerate}

\subsection{Objetivo Específico 3: Fase de Implementación del Prototipo}
\subsubsection{Actividades y Tareas}
\begin{enumerate}
	\item Construcción y Poblamiento del Motor de Datos:
	\begin{enumerate}
		\item Escribir los scripts DDL (Data Definition Language) para la creación de esquemas, implementando sentencias \texttt{CREATE TABLE} junto con sus restricciones de integridad (\texttt{PRIMARY KEY}, \texttt{FOREIGN KEY}, \texttt{NOT NULL}).
		\item Poblar las estructuras relacionales con un conjunto de datos de prueba (mock data) representativo mediante sentencias \texttt{INSERT INTO}, preparando el terreno para las pruebas funcionales.
	\end{enumerate}
	\item Programación de la Capa de Aplicación y Presentación:
	\begin{enumerate}
		\item Desarrollar el frontend utilizando semántica HTML5 y reglas de presentación CSS3, garantizando una correcta adaptabilidad en la visualización.
		\item Codificar la lógica de negocio en el backend mediante PHP, estableciendo conectores seguros (por ejemplo, PDO o conectores nativos para PostgreSQL/MySQL) hacia el SGBD.
		\item Implementar las rutinas transaccionales (CRUD) para todas las entidades y desarrollar funciones lógicas o \textit{triggers} en la base de datos que automaticen la integridad referencial profunda.
	\end{enumerate}
\end{enumerate}

\subsection{Objetivo Específico 4: Fase de Validación y Pruebas (Simulación)}
\subsubsection{Actividades y Tareas}
\begin{enumerate}
	\item Ejecución de Pruebas Unitarias y de Integración:
	\begin{enumerate}
		\item Diseñar casos de prueba aislados para validar las funciones críticas del backend PHP y el comportamiento de los \textit{triggers} en la base de datos.
		\item Comprobar la correcta intercomunicación de datos (integración) entre los formularios web, el servidor de aplicaciones y las tablas relacionales.
	\end{enumerate}
	\item Validación Estructural y de Seguridad del Prototipo:
	\begin{enumerate}
		\item Evaluar el rendimiento del motor de base de datos ejecutando consultas predefinidas y verificando el impacto en los tiempos de respuesta tras la creación de índices.
		\item Ejecutar pruebas de seguridad perimetral sobre el prototipo, orientadas a mitigar vulnerabilidades comunes como la Inyección SQL (SQLi) y el Cross-Site Scripting (XSS), consolidando la validación del sistema sin requerir un paso a producción.
	\end{enumerate}
\end{enumerate}

\section[Equipo de Trabajo]{Equipo de Desarrollo del Proyecto y Asignación de Responsabilidades}\label{equipo}

La estructura del equipo asume una distribución funcional basada en las fases CDI, delegando responsabilidades específicas para la conceptualización y construcción del prototipo tecnológico[cite: 1]. Esta configuración sugiere un trabajo interdisciplinario donde cada rol nutre las distintas etapas de modelado e implementación.

\noindent \textbf{1. Javier Rojas}
\begin{itemize}
	\item \textbf{Rol Principal:} Jefe de Proyecto \& Desarrollador Full-Stack Senior.
	\item \textbf{Responsabilidades Clave:}
	\begin{itemize}
		\item Coordinar los \textit{Sprints} y facilitar las ceremonias del marco metodológico Scrum.
		\item Supervisar la arquitectura del software, garantizando la viabilidad técnica entre el diseño y la implementación.
		\item Construir componentes complejos de la lógica de negocio en el backend (Objetivo 3).
		\item Guiar la transición entre las fases de implementación y validación, asegurando la consistencia del prototipo (Objetivo 4).
	\end{itemize}
\end{itemize}
\vspace{0.5em}

\noindent \textbf{2. Carolina Silva}
\begin{itemize}
	\item \textbf{Rol Principal:} Analista de Sistemas \& Diseñadora de Bases de Datos.
	\item \textbf{Responsabilidades Clave:}
	\begin{itemize}
		\item Dirigir la Fase de Concepción, diagnosticando requerimientos y extrayendo las entidades del dominio (Objetivo 1).
		\item Liderar el Modelado Conceptual y Lógico, proyectando el Modelo Entidad-Relación (MER) y ejecutando la normalización estructural hasta la 3FN (Objetivo 2).
		\item Elaborar la documentación técnica vinculada a la viabilidad y planificación inicial.
	\end{itemize}
\end{itemize}
\vspace{0.5em}

\noindent \textbf{3. Martín Pizarro}
\begin{itemize}
	\item \textbf{Rol Principal:} Desarrollador Backend \& Administrador de Base de Datos (DBA) Junior.
	\item \textbf{Responsabilidades Clave:}
	\begin{itemize}
		\item Materializar la estructura de datos mediante la escritura de scripts DDL y el poblamiento inicial con \textit{mock data} (Objetivo 3).
		\item Codificar las rutinas transaccionales (CRUD) en PHP y establecer los conectores persistentes hacia el SGBD (Objetivo 3).
		\item Implementar \textit{triggers} en el motor de base de datos para automatizar la integridad referencial.
	\end{itemize}
\end{itemize}
\vspace{0.5em}

\noindent \textbf{4. Sofía Castro}
\begin{itemize}
	\item \textbf{Rol Principal:} Desarrolladora Full-Stack \& Diseñadora UI/UX.
	\item \textbf{Responsabilidades Clave:}
	\begin{itemize}
		\item Esquematizar la experiencia de usuario mediante wireframes y mockups diferenciados por perfiles de acceso (Objetivo 2).
		\item Programar la capa de presentación utilizando semántica HTML5 y reglas de diseño CSS3 (Objetivo 3).
		\item Acoplar la interfaz gráfica con las rutinas del backend, evaluando la fluidez de interacción en el prototipo funcional.
	\end{itemize}
\end{itemize}
\vspace{0.5em}

\noindent \textbf{5. Diego Valdés}
\begin{itemize}
	\item \textbf{Rol Principal:} Desarrollador Full-Stack \& Soporte de Entornos.
	\item \textbf{Responsabilidades Clave:}
	\begin{itemize}
		\item Configurar el entorno de desarrollo controlado, habilitando el servidor web y los intérpretes necesarios para el sistema.
		\item Brindar apoyo transversal en la implementación de módulos CRUD y la escritura de funciones complementarias (Objetivo 3).
		\item Ejecutar evaluaciones estructurales preliminares sobre el motor de base de datos (Objetivo 4).
	\end{itemize}
\end{itemize}
\vspace{0.5em}

\noindent \textbf{6. Valentina Herrera}
\begin{itemize}
	\item \textbf{Rol Principal:} Encargada de Calidad (QA) \& Documentación Técnica.
	\item \textbf{Responsabilidades Clave:}
	\begin{itemize}
		\item Estructurar y ejecutar la Fase de Validación, diseñando casos de prueba unitaria y de integración para el prototipo (Objetivo 4).
		\item Realizar auditorías de seguridad perimetral (mitigación de SQLi y XSS) en un entorno simulado, sin pasar a producción.
		\item Consolidar la documentación técnica del proyecto, verificando la consistencia semántica y el apego a las normas académicas vigentes.
	\end{itemize}
\end{itemize}

\subsection{Responsabilidades Compartidas por Todo el Equipo}
\begin{itemize}
	\item \textbf{Comprensión de Requisitos:} Todos los miembros participarán en la discusión y comprensión de los requisitos del cliente.
	\item \textbf{Revisiones de Código:} Se realizarán revisiones cruzadas de código para mejorar la calidad y compartir conocimiento.
	\item \textbf{Pruebas:} Aunque Valentina lidera QA, todos los desarrolladores son responsables de las pruebas unitarias de su código y colaborarán en la identificación y corrección de errores.
	\item \textbf{Reuniones de Equipo:} Participación activa en reuniones periódicas de seguimiento, planificación y resolución de problemas, siguiendo la metodología Scrum.
	\item \textbf{Gestión del Conocimiento:} Compartir aprendizajes y soluciones encontradas durante el desarrollo.
	\item \textbf{Adaptabilidad:} Estar dispuestos a asumir tareas fuera de su rol principal si el proyecto lo requiere, dada la dinámica de un equipo pequeño y un proyecto formativo.
\end{itemize}




	\section{Planificación Gantt del Proyecto}
	
	La distribución cronológica presentada en la Figura \ref{fig:gantt_semanal} establece una secuencia trazable desde la abstracción de requerimientos hasta la validación empírica del prototipo computacional. El margen de tiempo semanal permite gestionar con precisión la concurrencia de tareas; por ejemplo, la estructuración del Modelo Entidad-Relación (MER) avanza en paralelo con la elaboración de la documentación técnica. Esta arquitectura de planificación podría facilitar una evaluación continua del rendimiento del equipo, mitigando el riesgo de desviaciones durante la fase de codificación.
	
	Los meses de enero y febrero de 2026 han sido excluidos deliberadamente del flujo continuo, considerándose periodos de inactividad académica.
	
	En la Tabla \ref{t:precedencia_wbs}  muestra las precedencias de cada tarea de esta carta Gantt.
\begin{landscape}	
	\begin{figure}[htb]
		\centering
		\resizebox{\linewidth}{!}{
			\begin{ganttchart}[
				x unit=0.25cm,
				y unit title=0.5cm,
				y unit chart=0.55cm,
				vgrid,
				hgrid,
				title height=1,
				bar/.style={fill=blue!70, draw=blue!90},
				bar incomplete/.append style={fill=blue!10, draw=blue!40},
				progress label text={#1\%},
				group incomplete/.append style={draw=black, fill=black!20},
				group/.append style={draw=black, fill=black},
				link/.style={-latex, line width=1pt, blue!80},
				link label font=\scriptsize\bfseries
				]{1}{60}
				
				% Títulos de años y meses (4 semanas por mes = 60 semanas total)
				\gantttitle{Año 2025}{40} \gantttitle{Año 2026 (Excluye Ene-Feb)}{20} \\
				\gantttitle{Mar}{4} \gantttitle{Abr}{4} \gantttitle{May}{4} \gantttitle{Jun}{4} \gantttitle{Jul}{4} \gantttitle{Ago}{4} \gantttitle{Sep}{4} \gantttitle{Oct}{4} \gantttitle{Nov}{4} \gantttitle{Dic}{4} 
				\gantttitle{Mar}{4} \gantttitle{Abr}{4} \gantttitle{May}{4} \gantttitle{Jun}{4} \gantttitle{Jul}{4} \\
				% Escala Semanal
				\gantttitlelist{1,...,60}{1} \\
				
				% ==========================================
				% OBJETIVO 1: CONCEBIR (100%)
				% ==========================================
				\ganttgroup[progress=100]{O.E.1: Fase de Concepción}{1}{16} \\
				\ganttgroup[progress=100]{Actividad 1.1: Diagnóstico y Entidades}{1}{8} \\
				\ganttbar[progress=100, name=T111]{T1.1.1: Levantamiento de requerimientos}{1}{4} \\
				\ganttbar[progress=100, name=T112]{T1.1.2: Evaluación de factibilidad técnica}{5}{8} \\
				\ganttgroup[progress=100]{Actividad 1.2: Consolidación Documental}{9}{16} \\
				\ganttbar[progress=100, name=T121]{T1.2.1: Estructuración informe técnico}{9}{12} \\
				\ganttbar[progress=100, name=T122]{T1.2.2: Planificación Gantt (Simultánea)}{9}{12} \\
				\ganttbar[progress=100, name=T123]{T1.2.3: Validación con cliente académico}{13}{16} \\
				
				% ==========================================
				% OBJETIVO 2: DISEÑAR (50%)
				% ==========================================
				\ganttgroup[progress=50]{O.E.2: Fase de Diseño}{17}{32} \\
				\ganttgroup[progress=75]{Actividad 2.1: Modelado de Datos}{17}{24} \\
				\ganttbar[progress=100, name=T211]{T2.1.1: Proyección del MER y atributos}{17}{20} \\
				\ganttbar[progress=50, name=T212]{T2.1.2: Transformación relacional y 3FN}{21}{24} \\
				\ganttgroup[progress=25]{Actividad 2.2: Arquitectura de Interfaces}{21}{32} \\
				\ganttbar[progress=50, name=T221]{T2.2.1: Wireframes de bajo nivel (Simultánea)}{21}{24} \\
				\ganttbar[progress=0, name=T222]{T2.2.2: Mockups alta fidelidad y UX}{25}{32} \\
				
				% ==========================================
				% OBJETIVO 3: IMPLEMENTAR (0%)
				% ==========================================
				\ganttgroup[progress=0]{O.E.3: Fase de Implementación}{33}{48} \\
				\ganttgroup[progress=0]{Actividad 3.1: Motor de Datos}{33}{40} \\
				\ganttbar[progress=0, name=T311]{T3.1.1: Scripts DDL y despliegue SGBD}{33}{36} \\
				\ganttbar[progress=0, name=T312]{T3.1.2: Poblamiento con Mock Data}{37}{40} \\
				% Salto lógico: semanas 41 a 48 corresponden a Marzo-Abril 2026 tras el receso.
				\ganttgroup[progress=0]{Actividad 3.2: Sistema Web y Lógica}{41}{48} \\
				\ganttbar[progress=0, name=T321]{T3.2.1: Programación Frontend (HTML5/CSS3)}{41}{44} \\
				\ganttbar[progress=0, name=T322]{T3.2.2: Conectores y CRUD Backend (PHP)}{45}{48} \\
				\ganttbar[progress=0, name=T323]{T3.2.3: Triggers de integridad (Simultánea)}{45}{48} \\
				
				% ==========================================
				% OBJETIVO 4: VALIDAR (0%)
				% ==========================================
				\ganttgroup[progress=0]{O.E.4: Fase de Validación}{49}{60} \\
				\ganttgroup[progress=0]{Actividad 4.1: Pruebas Unitarias e Integración}{49}{54} \\
				\ganttbar[progress=0, name=T411]{T4.1.1: Casos de prueba backend}{49}{51} \\
				\ganttbar[progress=0, name=T412]{T4.1.2: Verificación de flujos transaccionales}{52}{54} \\
				\ganttgroup[progress=0]{Actividad 4.2: Evaluación Estructural}{55}{60} \\
				\ganttbar[progress=0, name=T421]{T4.2.1: Análisis de rendimiento en consultas}{55}{57} \\
				\ganttbar[progress=0, name=T422]{T4.2.2: Pruebas mitigación SQLi/XSS (Simultánea)}{55}{57} \\
				\ganttbar[progress=0, name=T423]{T4.2.3: Consolidación final del prototipo}{58}{60} \\
				
				% ==========================================
				% DEPENDENCIAS LÓGICAS (ENLACES)
				% ==========================================
				\ganttlink{T111}{T112}
				\ganttlink{T112}{T121}
				\ganttlink{T121}{T123}
				\ganttlink{T123}{T211}
				\ganttlink{T211}{T212}
				\ganttlink{T222}{T311}
				\ganttlink{T311}{T312}
				\ganttlink{T312}{T321}
				\ganttlink{T321}{T322}
				\ganttlink{T322}{T411}
				\ganttlink{T411}{T412}
				\ganttlink{T412}{T421}
				\ganttlink{T421}{T423}
				
			\end{ganttchart}
		}
		\caption{Desglose semanal de la Carta Gantt del proyecto, visibilizando objetivos, actividades y tareas. Se evidencia el estado actual de progreso y el paralelismo operativo aplicado.}
		\label{fig:gantt_semanal}
	\end{figure}
	
\end{landscape}




\begin{table}[htb]
	\centering	
	\begin{tabular}{|c|c|c|c|c|}\hline
		\textbf{WBS}&\textbf{Duración} &\textbf{Después de} & \textbf{Simultanea a} & \textbf{Antes de}\\\hline
		T1.1.1 & 4 & - & - & T1.1.2 \\\hline
		T1.1.2 & 4 & T1.1.1 & - & T1.2.1, T1.2.2 \\\hline
		T1.2.1 & 4 & T1.1.2 & T1.2.2 & T1.2.3 \\\hline
		T1.2.2 & 4 & T1.1.2 & T1.2.1 & T1.2.3 \\\hline
		T1.2.3 & 4 & T1.2.1, T1.2.2 & - & T2.1.1 \\\hline
		T2.1.1 & 4 & T1.2.3 & - & T2.1.2, T2.2.1 \\\hline
		T2.1.2 & 4 & T2.1.1 & T2.2.1 & T2.2.2 \\\hline
		T2.2.1 & 4 & T2.1.1 & T2.1.2 & T2.2.2 \\\hline
		T2.2.2 & 8 & T2.1.2, T2.2.1 & - & T3.1.1 \\\hline
		T3.1.1 & 4 & T2.2.2 & - & T3.1.2 \\\hline
		T3.1.2 & 4 & T3.1.1 & - & T3.2.1 \\\hline
		T3.2.1 & 4 & T3.1.2 & - & T3.2.2, T3.2.3 \\\hline
		T3.2.2 & 4 & T3.2.1 & T3.2.3 & T4.1.1 \\\hline
		T3.2.3 & 4 & T3.2.1 & T3.2.2 & T4.1.1 \\\hline
		T4.1.1 & 3 & T3.2.2, T3.2.3 & - & T4.1.2 \\\hline
		T4.1.2 & 3 & T4.1.1 & - & T4.2.1, T4.2.2 \\\hline
		T4.2.1 & 3 & T4.1.2 & T4.2.2 & T4.2.3 \\\hline
		T4.2.2 & 3 & T4.1.2 & T4.2.1 & T4.2.3 \\\hline
		T4.2.3 & 3 & T4.2.1, T4.2.2 & - & - \\\hline
	\end{tabular}
	\caption{Matriz de precedencias estructurada según nomenclatura WBS, evidenciando dependencias y tareas concurrentes.}
	\label{t:precedencia_wbs}	
\end{table}
 

 
\section[Justificación y Aporte]{Justificación del Proyecto, Utilidad y Aporte para la Universidad}

\subsection{Justificación del Proyecto}
La gestión eficiente de la información académica y de investigación es un pilar fundamental para el funcionamiento y la excelencia de cualquier institución de educación superior. En el contexto actual, donde las universidades manejan un volumen creciente y complejo de datos relacionados con su personal docente, alumnado, proyectos de investigación, programas de postgrado y estructura departamental, la necesidad de contar con herramientas tecnológicas robustas y centralizadas se vuelve imperativa.

Actualmente, es común que esta información se encuentre dispersa en múltiples sistemas aislados, hojas de cálculo o incluso en formatos manuales, lo que genera ineficiencias, dificulta la toma de decisiones informadas, incrementa el riesgo de inconsistencias y limita la capacidad de análisis integral del quehacer universitario. El presente proyecto se justifica por la necesidad de superar estas limitaciones mediante el desarrollo de un sistema de base de datos web diseñado a medida, que permita una administración integrada, ágil y precisa de todos estos datos críticos.

\subsection{Utilidad del Sistema para la Universidad}

El sistema de base de datos web propuesto ofrecerá una utilidad transversal a diversas áreas y niveles de la Universidad:
\begin{enumerate}
    \item \textbf{Centralización e Integración de la Información}: Se consolidarán los datos de profesores, alumnos, proyectos, departamentos, programas de postgrado y asignaturas en una única plataforma. Esto eliminará redundancias, mejorará la consistencia de los datos y proporcionará una visión unificada de las actividades académicas y de investigación.
\item \textbf{Optimización de Procesos Administrativos}: Se simplificarán y agilizarán numerosas tareas administrativas, tales como:
\begin{itemize}
    \item La asignación de profesores a proyectos y la gestión de sus roles (investigador principal, investigador).
\item El seguimiento de la participación de alumnos de postgrado en proyectos como ayudantes de investigación y la asignación de sus supervisores.
\item La administración de la carga docente y la afiliación de profesores a departamentos con sus respectivos porcentajes de dedicación.
\item La gestión de la información de los programas de postgrado y la adscripción de los alumnos.
La asignación de directores de departamento y asesores de alumnos.
\end{itemize}
\item \textbf{Facilitación del Acceso a la Información}: Permitirá a los usuarios autorizados (desde personal administrativo hasta directivos y los propios académicos y alumnos, según sus perfiles) acceder de manera rápida y segura a la información que necesitan, mejorando la transparencia y la capacidad de respuesta.
\item \textbf{Soporte para la Generación de Informes y Estadísticas}: Al contar con datos estructurados y centralizados, el sistema facilitará la generación de informes detallados y estadísticas relevantes para la evaluación del desempeño, la planificación estratégica y la acreditación institucional. Por ejemplo, se podrá conocer fácilmente la productividad investigadora, la carga académica por departamento o la tasa de participación de alumnos en proyectos.
\end{enumerate}

\subsection{Aporte del Sistema a la Universidad}
La implementación de este sistema representará un aporte significativo para la Universidad en múltiples dimensiones:
\begin{enumerate}
    \item \textbf{Mejora de la Eficiencia Operativa}:
\begin{itemize}
    \item Reducción del tiempo dedicado a la búsqueda y consolidación manual de información.
\item Automatización de procesos clave, disminuyendo la carga de trabajo administrativo y la probabilidad de errores humanos.
\item Optimización en la asignación de recursos humanos y financieros en proyectos y departamentos.
\end{itemize}
\item \textbf{Fortalecimiento de la Toma de Decisiones Estratégicas}:
\begin{itemize}
\item Proporcionará a la gerencia y a los responsables de unidades académicas información fiable y oportuna para la planificación a corto, mediano y largo plazo.
\item Permitirá identificar fortalezas, debilidades, oportunidades y amenazas en el ámbito académico y de investigación con base en evidencia concreta.
 \item Facilitará la evaluación del impacto de programas y políticas institucionales.
 \end{itemize}
 \item \textbf{Impulso a la Gestión de la Investigación}:
\begin{itemize}
\item Mejorará el seguimiento de los proyectos de investigación: sus participantes, financiamiento, plazos y supervisión.
\item Facilitará la identificación de sinergias y posibles colaboraciones entre investigadores y departamentos.
\item Contribuirá a la visibilidad de la producción científica de la institución.
\end{itemize}
 \item \textbf{Mejora en la Gestión y Apoyo Académico}:
\begin{itemize}
    \item Optimizará la administración de los programas de postgrado y el seguimiento de sus estudiantes.
\item Facilitará la labor de supervisión y asesoría a los alumnos.
\item Permitirá una gestión más clara de las responsabilidades docentes y de investigación del cuerpo académico.
\end{itemize}

\item \textbf{Incremento de la Transparencia y Rendición de Cuentas}:
\begin{itemize}
\item Proporcionará un registro claro y auditable de las actividades y asignaciones.
\item Facilitará la preparación de informes para organismos de acreditación, entidades financiadoras y la comunidad universitaria en general.
\end{itemize}

\item\textbf{Modernización Tecnológica y Base para el Futuro}:

\begin{itemize}
\item Representa un paso adelante en la modernización de la infraestructura tecnológica de la Universidad.
\item Un sistema bien diseñado y escalable podrá adaptarse a futuras necesidades e integrarse con otras plataformas institucionales, consolidando el ecosistema digital de la Universidad.
\end{itemize}
\end{enumerate}
En resumen, este sistema no solo resolverá problemas operativos existentes, sino que también se convertirá en una herramienta estratégica que potenciará la capacidad de gestión, la calidad de la investigación y la formación académica, y la eficiencia general de la Universidad, contribuyendo directamente al cumplimiento de su misión y visión institucional.

\section{Análisis de Viabilidad del Proyecto}
La implementación del sistema propuesto para la Universidad se considera viable tras analizar los siguientes aspectos:

\subsection{Viabilidad Técnica}
El proyecto es técnicamente viable por las siguientes razones:
\begin{itemize}
    \item \textbf{Tecnologías Maduras y Disponibles}: Las tecnologías propuestas para el desarrollo (HTML, CSS, JavaScript para el frontend; PHP para el backend; un Sistema Gestor de Base de Datos SQL como MySQL o PostgreSQL) son estándares en la industria, maduras, ampliamente documentadas y con una gran comunidad de soporte. Existen abundantes herramientas de desarrollo, muchas de ellas de código abierto, que facilitan su implementación.
\item \textbf{Complejidad Manejable}: Si bien el sistema integra diversas entidades y relaciones, la lógica de negocio (operaciones CRUD, gestión de roles, reglas de integridad) es abordable con las tecnologías mencionadas. El modelado relacional y la normalización (hasta 3FN) asegurarán una base de datos robusta y eficiente. La implementación de funciones y triggers (PL/SQL o similar, según el SGBD) está dentro de las capacidades estándar de los SGBD modernos.
\item \textbf{Recursos Humanos Calificados}: Se asume la disponibilidad de un equipo de desarrollo con experiencia en análisis de sistemas, diseño de bases de datos, desarrollo web full-stack (PHP, HTML, CSS) y gestión de proyectos.
\item \textbf{Infraestructura Requerida}: La infraestructura necesaria (servidores web y de base de datos) puede ser implementada en las instalaciones de la Universidad o mediante servicios de cloud computing, ofreciendo flexibilidad y escalabilidad. Se considera que la Universidad cuenta con la capacidad para alojar o contratar estos servicios.
\item \textbf{Escalabilidad y Rendimiento}: El diseño propuesto, incluyendo una correcta normalización, indexación y la posibilidad de optimizar consultas, permitirá que el sistema maneje el volumen de datos y la carga de usuarios esperada en una institución universitaria. El sistema podrá escalar horizontal o verticalmente según las necesidades futuras.

\end{itemize}
\subsection{Viabilidad Económica}
Se estima que el proyecto es económicamente viable, considerando:

\subsubsection{Costos de Desarrollo}
\begin{itemize}
    \item \textbf{Personal}: Principal componente de costo, incluyendo analistas, diseñadores, desarrolladores (frontend y backend), especialistas en bases de datos y testers durante los 15 meses.
\item \textbf{Software}: Potencialmente, costos de licencias para el SGBD o herramientas de desarrollo especializadas, aunque se priorizará el uso de software de código abierto o con licenciamiento favorable para instituciones educativas.
\item \textbf{Hardware/Infraestructura}: Costos iniciales de adquisición o configuración de servidores, o costos recurrentes si se opta por servicios en la nube.
\end{itemize}
\subsubsection{Costos Operativos (Post-Implementación)}
\begin{itemize}
    \item Mantenimiento del software y la base de datos.
\item Soporte técnico a usuarios.
\item Costos continuos de infraestructura (energía, conectividad, cloud si aplica).
\item Capacitación continua.
\end{itemize}
\subsubsection{Retorno de la Inversión (ROI) y Beneficios} Aunque muchos beneficios son cualitativos, su impacto económico es considerable:
\begin{itemize}
    \item \textbf{Reducción de Costos Directos}: Disminución de horas-hombre dedicadas a tareas manuales de gestión de datos, generación de informes y resolución de inconsistencias.
\item \textbf{Eficiencia Mejorada}: Optimización de procesos que liberan tiempo de personal cualificado para tareas de mayor valor añadido.
\item \textbf{Mejora en la Toma de Decisiones}: Decisiones basadas en datos precisos y oportunos pueden llevar a una asignación más eficiente de recursos, optimización de programas académicos y mejora en la captación de fondos de investigación.
\item \textbf{Cumplimiento y Acreditación}: Un sistema robusto facilita el cumplimiento de normativas y los procesos de acreditación, evitando posibles sanciones o asegurando financiamiento.
\item \textbf{Reducción de Errores}: La automatización y validaciones disminuyen la incidencia de errores costosos.
\end{itemize}
La inversión inicial se justifica por los ahorros a mediano y largo plazo, la mejora en la calidad de la gestión y el soporte a las funciones críticas de la Universidad. Se recomienda realizar un análisis costo-beneficio más detallado una vez afinados los requerimientos técnicos específicos y la selección de plataformas.

\subsection{Viabilidad Operativa}

El sistema se considera operativamente viable, ya que:
\begin{itemize}
    \item \textbf{Alineación con Necesidades del Usuario}: El sistema se diseñará basándose en los requisitos específicos proporcionados, buscando resolver problemáticas concretas de la gestión académica y de investigación. Las interfaces diferenciadas para usuarios y administradores con distintos niveles de funcionalidad CRUD están pensadas para adaptarse a sus roles.
\item \textbf{Aceptación por parte de los Usuarios}: Se mitigarán resistencias al cambio mediante:
\begin{itemize}
\item \textbf{Capacitación}: Se planificarán sesiones de capacitación para los diferentes perfiles de usuario.
\item \textbf{Diseño Intuitivo}: Se priorizará una interfaz de usuario clara, amigable y fácil de usar.
\item \textbf{Participación}: Involucrar a representantes de los usuarios finales durante las fases de diseño y pruebas de aceptación (UAT) para asegurar que el sistema cumpla con sus expectativas y facilite su trabajo diario.
\end{itemize}
\item \textbf{Integración en Flujos de Trabajo}: El sistema está concebido para integrarse y optimizar los flujos de trabajo existentes. Es posible que se requiera una adaptación de ciertos procesos manuales para aprovechar al máximo las capacidades del nuevo sistema, lo cual se gestionará como parte del proceso de implementación.
\item \textbf{Soporte y Mantenimiento}: Se deberá establecer un plan de soporte técnico continuo (interno o externo) para resolver incidencias, realizar mantenimiento preventivo y aplicar futuras actualizaciones o mejoras. La documentación técnica y de usuario generada facilitará estas labores.
\item \textbf{Seguridad de la Información}: Se implementarán medidas de seguridad adecuadas para proteger la integridad, confidencialidad y disponibilidad de los datos, incluyendo gestión de accesos basada en roles, protección contra amenazas comunes y políticas de respaldo.
\end{itemize}


\subsection{Viabilidad de Plazos (Cronograma – 15 Meses)}

Un plazo total de 15 meses de trabajo efectivo se considera factible para la realización de este proyecto, dada su envergadura y la adopción de la metodología Scrum. Este marco de trabajo ágil permitirá entregas incrementales y una adaptación continua a lo largo del desarrollo.

La planificación temporal se estructura de la siguiente manera, considerando los períodos de receso académico:
\begin{itemize}
    \item \textbf{Inicio del Proyecto:} Septiembre 2025.
    \item \textbf{Duración Total de Trabajo Efectivo:} 15 meses.
    \item \textbf{Meses No Laborables (Receso Académico):} Febrero 2026 y Agosto 2026.
    \item \textbf{Finalización Estimada del Proyecto:} Enero 2027.
\end{itemize}

El desarrollo se organizará en Sprints (ciclos de trabajo cortos, por ejemplo, de 4 semanas). Los objetivos específicos se distribuirán en bloques de trabajo de la siguiente forma:

\begin{itemize}
    \item \textbf{Bloque 1: Objetivo Específico 1 (Formulación y Modelado Conceptual).}
    \begin{itemize}
        \item \textbf{Período de Trabajo:} Septiembre 2025 – Enero 2026 (5 meses de trabajo).
        \item \textbf{Sprints Estimados:} Aproximadamente 5 Sprints.
        \item \textbf{Enfoque Principal:} Dedicado a la elaboración del informe técnico inicial (incluyendo análisis de viabilidad, justificación, etc.), análisis exhaustivo de requisitos, y el diseño detallado del Modelo Entidad-Relación (MER) y el Modelo Conceptual de la base de datos.
    \end{itemize}

    \item \textbf{Bloque 2: Objetivo Específico 2 (Diseño e Implementación de la Base de Datos).}
    \begin{itemize}
        \item \textbf{Período de Trabajo:} Marzo 2026 – Julio 2026 (5 meses de trabajo, posterior al receso de Febrero 2026).
        \item \textbf{Sprints Estimados:} Aproximadamente 5 Sprints.
        \item \textbf{Enfoque Principal:} Centrado en la transformación del modelo conceptual al modelo relacional, aplicación de técnicas de normalización (hasta 3FN), escritura de los scripts DDL para la creación de tablas y restricciones, carga de datos de prueba, y la implementación de índices y vistas predefinidas para optimizar consultas.
    \end{itemize}

    \item \textbf{Bloque 3: Objetivos Específicos 3 y 4 (Diseño y Desarrollo del Sistema Web, Integración con la Base de Datos, Procesos CRUD, Pruebas y Validación).}
    \begin{itemize}
        \item \textbf{Período de Trabajo:} Septiembre 2026 – Enero 2027 (5 meses de trabajo, posterior al receso de Agosto 2026).
        \item \textbf{Sprints Estimados:} Aproximadamente 5 Sprints.
        \item \textbf{Enfoque Principal:} Esta fase comprenderá el diseño de maquetas del sistema web, el desarrollo frontend (HTML, CSS) y backend (PHP), la conexión con la base de datos implementada, y la programación de las funcionalidades CRUD. De forma paralela e integrada en los Sprints, se llevará a cabo el Objetivo 4, realizando pruebas unitarias, de integración, de sistema, y las pruebas de aceptación del usuario (UAT). Se finalizará con la corrección de errores, la elaboración de la documentación final y la preparación para el despliegue.
    \end{itemize}
\end{itemize}

\textbf{Consideraciones para cumplir el plazo:}
El cumplimiento de este cronograma dependerá de varios factores críticos:
\begin{itemize}
    \item \textbf{Gestión de Riesgos Eficaz:} Identificación temprana y mitigación proactiva de posibles riesgos, como cambios no controlados en el alcance (scope creep), retrasos en la toma de decisiones o validaciones por parte del ``cliente'' (en este caso, el profesorado del curso), o la aparición de problemas técnicos imprevistos.
    \item \textbf{Dedicación y Compromiso del Equipo:} Es fundamental contar con la dedicación adecuada del equipo de desarrollo estudiantil y un compromiso con los objetivos de cada Sprint y las entregas planificadas.
    \item \textbf{Aplicación Rigurosa de Scrum:} La correcta implementación de las ceremonias (Sprint Planning, Daily Scrum, Sprint Review, Sprint Retrospective) y artefactos (Product Backlog, Sprint Backlog, Incremento) de Scrum facilitará la transparencia, la inspección y la adaptación continua, elementos clave para el manejo eficiente del tiempo.
    \item \textbf{Comunicación Constante y Fluida:} Mantener una comunicación efectiva dentro del equipo de desarrollo y con el ``cliente'' para asegurar la alineación y la rápida resolución de dudas o problemas que puedan surgir.
\end{itemize}

En conclusión, el plazo de 15 meses de trabajo efectivo, distribuido según el calendario propuesto y gestionado mediante la metodología ágil Scrum, ofrece un marco temporal realista y estructurado para la finalización exitosa del proyecto por parte del equipo de estudiantes.





\section[Metodología de Desarrollo]{Metodología de Desarrollo Propuesta: Scrum}
Para un equipo de 6 estudiantes de Ingeniería Civil en Informática que están comenzando a conocer las metodologías de desarrollo de software y que tienen un proyecto de 15 meses, la metodología más apropiada sería \textbf{Scrum}.

\subsection*{Justificación de la Elección de Scrum}

Scrum es un marco de trabajo ágil que se adapta muy bien a las características de este proyecto y del equipo por las siguientes razones:

\begin{enumerate}
    \item \textbf{Desarrollo Iterativo e Incremental:}
    Scrum divide el proyecto en ciclos cortos llamados \textit{Sprints} (generalmente de 2 a 4 semanas). Al final de cada Sprint, el equipo entrega un incremento funcional del producto. Esto es ideal para un proyecto de 15 meses, ya que permite al equipo mostrar progreso tangible de forma regular, obtener retroalimentación temprana del ``cliente'' (profesor del curso) y ajustar el rumbo si es necesario. Para estudiantes, ver resultados funcionales de forma periódica es altamente motivador.

    \item \textbf{Adaptabilidad y Flexibilidad:}
    Aunque los requisitos iniciales están definidos, es probable que surjan nuevos entendimientos o se necesiten ajustes, especialmente en la interfaz de usuario y la experiencia de usuario (UI/UX). Scrum permite la adaptación a cambios entre Sprints, incorporando nuevos requisitos o modificaciones en el \textit{Product Backlog}.

    \item \textbf{Estructura y Claridad para Equipos Nuevos:}
    Para estudiantes que recién se inician en metodologías, Scrum ofrece un marco con roles, eventos y artefactos definidos que proporcionan una estructura clara:
    \begin{itemize}
        \item \textit{Roles Adaptados:} Para el contexto estudiantil, estos roles pueden ser:
        \begin{itemize}
            \item \textbf{Product Owner (Dueño del Producto):} Rol asumido idealmente por el ``cliente'' (profesor) o, de forma interna y en representación del cliente, por un miembro del equipo (ej. el Jefe de Proyecto, Javier Rojas) encargado de definir y priorizar los ítems del Product Backlog.
            \item \textbf{Scrum Master (Facilitador):} Rol que podría ser asumido por el Jefe de Proyecto (Javier Rojas), enfocándose en asegurar que el equipo siga las prácticas de Scrum, facilitar las reuniones y eliminar impedimentos.
            \item \textbf{Development Team (Equipo de Desarrollo):} Los 6 estudiantes, como un equipo auto-organizado y multidisciplinario, responsable de entregar el incremento del producto.
        \end{itemize}
        \item \textit{Eventos (Reuniones):}
        \begin{itemize}
            \item \textit{Sprint Planning:} Planificación del trabajo a realizar en el Sprint.
            \item \textit{Daily Scrum:} Reunión diaria corta para sincronizar al equipo.
            \item \textit{Sprint Review:} Demostración del incremento del producto y recolección de feedback.
            \item \textit{Sprint Retrospective:} Reflexión sobre el Sprint anterior para identificar mejoras en el proceso.
        \end{itemize}
        Estas reuniones estructuradas ayudan a mantener el enfoque y la comunicación.
        \item \textit{Artefactos:}
        \begin{itemize}
            \item \textit{Product Backlog:} Lista priorizada de todas las funcionalidades y requisitos del proyecto.
            \item \textit{Sprint Backlog:} Conjunto de ítems del Product Backlog seleccionados para un Sprint, más el plan para entregarlos.
            \item \textit{Incremento del Producto:} La suma de todos los ítems del Product Backlog completados durante un Sprint y Sprints anteriores, que debe ser potencialmente entregable.
        \end{itemize}
    \end{itemize}
    Esta estructura es fundamental para guiar al equipo a lo largo de los 15 meses de proyecto.

    \item \textbf{Fomenta la Colaboración y Comunicación:}
    El tamaño del equipo (6 personas) es ideal para Scrum. Las reuniones como el Daily Scrum y la planificación y revisión conjuntas de los Sprints aseguran que todos los miembros estén alineados, compartan información y colaboren estrechamente para alcanzar el objetivo del Sprint.

    \item \textbf{Gestión de Riesgos Temprana:}
    La naturaleza iterativa y las revisiones constantes (Sprint Review, Sprint Retrospective) permiten identificar problemas, riesgos o malentendidos de forma temprana, en lugar de esperar hasta el final del proyecto. Esto es crucial en un proyecto de 15 meses, donde los problemas no detectados pueden escalar significativamente.

    \item \textbf{Enfoque en la Entrega de Valor:}
    Cada Sprint se enfoca en entregar las funcionalidades de mayor valor para el cliente (priorizadas en el Product Backlog), lo que asegura que el sistema evolucione de manera útil y relevante desde las primeras etapas.

    \item \textbf{Oportunidad de Aprendizaje Práctico:}
    Implementar Scrum proporcionará al equipo una valiosa experiencia práctica con una de las metodologías ágiles más utilizadas en la industria del software. Esto constituye un objetivo de aprendizaje importante en su formación como Ingenieros Civiles en Informática, preparándolos para entornos de trabajo reales.
\end{enumerate}

Aunque otras metodologías ágiles como Kanban podrían ser consideradas por su simplicidad visual y enfoque en el flujo, la estructura de Sprints de Scrum, junto con sus roles y ceremonias definidas, ofrece un andamiaje más robusto y guiado para un equipo de estudiantes que se enfrenta a un proyecto de duración considerable y que necesita aprender a gestionar su trabajo de forma sistemática, colaborativa y adaptativa. Les ayudará a mantener un ritmo sostenible y a asegurar entregas consistentes y de valor a lo largo de los 15 meses del proyecto.


\section[Planificación Temporal]{Planificación Temporal del Desarrollo del Proyecto (Metodología Scrum)}

El marco ágil Scrum estructura el avance en iteraciones, ajustando los incrementos al ciclo de vida del prototipo[cite: 1]. La calendarización distribuye el esfuerzo a lo largo de quince meses efectivos, omitiendo los recesos académicos y fragmentando el trabajo en bloques orientados a las fases CDI y su consecuente validación funcional.

\noindent\textbf{Consideraciones Metodológicas de la Planificación:}
\begin{itemize}
	\item \textbf{Inicio del Proyecto:} Septiembre 2025.
	\item \textbf{Duración Total:} 15 meses de trabajo efectivo (excluyendo recesos de Febrero y Agosto 2026).
	\item \textbf{Iteraciones:} Sprints de 4 semanas, subordinados al refinamiento continuo del \textit{Product Backlog}.
	\item \textbf{Delimitación del Alcance:} Las entregas finales se evalúan bajo métricas de prototipado, descartando actividades de despliegue operativo comercial.
\end{itemize}

\begin{landscape}
	{\small
		\begin{longtable}{>{\RaggedRight}p{2.8cm} >{\Centering}p{1.3cm} >{\RaggedRight}p{2cm} >{\Centering}p{1cm} >{\RaggedRight}p{1.8cm} >{\RaggedRight}p{5.5cm} >{\RaggedRight}p{5cm}}
			\caption{Planificación Temporal del Proyecto Orientado a Prototipo}\label{tab:planificacion}\\
			\toprule
			\textbf{Fase / Bloque de Objetivos} & \textbf{Mes de Trabajo} & \textbf{Mes y Año Calendario} & \textbf{Sprint \#} & \textbf{Duración Estimada} & \textbf{Foco Principal y Actividades Clave} & \textbf{Entregables Clave del Sprint} \\
			\midrule
			\endfirsthead
			\caption[]{Planificación Temporal del Proyecto (Continuación)}\\
			\toprule
			\textbf{Fase / Bloque de Objetivos} & \textbf{Mes de Trabajo} & \textbf{Mes y Año Calendario} & \textbf{Sprint \#} & \textbf{Duración Estimada} & \textbf{Foco Principal y Actividades Clave} & \textbf{Entregables Clave del Sprint} \\
			\midrule
			\endhead
			\midrule
			\multicolumn{7}{r}{{\footnotesize Continuará en la página siguiente...}} \\
			\midrule
			\endfoot
			\bottomrule
			\endlastfoot
			
			\textbf{Bloque 1: Fase de Concepción} (Obj. Específico 1) & 1 & Sep 2025 & 1 & 4 semanas & Configuración del entorno Scrum. Levantamiento de requerimientos funcionales y extracción de entidades del dominio académico. & Acta de constitución, \textit{Product Backlog} inicial, listado de requerimientos y entidades. \\
			\midrule
			\textbf{Bloque 1: Fase de Concepción} (Obj. Específico 1) & 2 & Oct 2025 & 2 & 4 semanas & Estructuración del informe técnico preliminar. Evaluación de factibilidad técnica para el desarrollo web. & Borrador de informe técnico, matriz de viabilidad. \\
			\midrule
			\textbf{Bloque 1: Fase de Concepción} (Obj. Específico 1) & 3 & Nov 2025 & 3 & 4 semanas & Definición exhaustiva de atributos y restricciones. Inicio del diseño del Modelo Entidad-Relación (MER). & Especificación de atributos, borrador inicial del MER. \\
			\midrule
			\textbf{Bloque 1: Fase de Concepción} (Obj. Específico 1) & 4 & Dic 2025 & 4 & 4 semanas & Refinamiento del diagrama conceptual. Consolidación de cardinalidades (1:N, M:N). & Modelo Entidad-Relación finalizado. \\
			\midrule
			\textbf{Bloque 1: Fase de Concepción} (Obj. Específico 1) & 5 & Ene 2026 & 5 & 4 semanas & Validación del modelo conceptual con contraparte académica. Revisión final del informe técnico. & Documento de formulación aprobado. \\
			\midrule
			\textit{Receso Académico} & -- & \textit{Feb 2026} & -- & -- & \textit{Suspensión temporal de actividades} & -- \\
			\midrule
			\textbf{Bloque 2: Fase de Diseño} (Obj. Específico 2) & 6 & Mar 2026 & 6 & 4 semanas & Transformación del MER al modelo relacional. Aplicación de normalización hasta Primera Forma Normal (1FN). & Modelo relacional preliminar. \\
			\midrule
			\textbf{Bloque 2: Fase de Diseño} (Obj. Específico 2) & 7 & Abr 2026 & 7 & 4 semanas & Culminación de normalización estructural (3FN). Elaboración de wireframes de bajo nivel para interfaces. & Esquemas en 3FN, wireframes de estructura base. \\
			\midrule
			\textbf{Bloque 2: Fase de Diseño} (Obj. Específico 2) & 8 & May 2026 & 8 & 4 semanas & Diseño de mockups de alta fidelidad diferenciando perfiles (administrador/limitado). Arquitectura de datos final. & Interfaces gráficas conceptuales, arquitectura de software. \\
			\midrule
			\textbf{Bloque 3: Fase de Implementación} (Obj. Específico 3) & 9 & Jun 2026 & 9 & 4 semanas & Escritura de scripts DDL. Creación física de tablas y llaves foráneas. Poblamiento con \textit{mock data}. & Base de datos estructurada y con datos de prueba. \\
			\midrule
			\textbf{Bloque 3: Fase de Implementación} (Obj. Específico 3) & 10 & Jul 2026 & 10 & 4 semanas & Desarrollo de maquetación HTML5 y reglas CSS3. Configuración del servidor de desarrollo controlado. & Vistas frontend estáticas operativas. \\
			\midrule
			\textit{Receso Académico} & -- & \textit{Ago 2026} & -- & -- & \textit{Suspensión temporal de actividades} & -- \\
			\midrule
			\textbf{Bloque 3: Fase de Implementación} (Obj. Específico 3) & 11 & Sep 2026 & 11 & 4 semanas & Programación de conectores PDO/PHP. Implementación de operaciones CRUD para entidades primarias (Profesores, Alumnos). & Módulos CRUD funcionales integrados a la BD. \\
			\midrule
			\textbf{Bloque 3: Fase de Implementación} (Obj. Específico 3) & 12 & Oct 2026 & 12 & 4 semanas & Desarrollo de lógica de negocio compleja (asignaciones a proyectos). Implementación de \textit{triggers} en el SGBD. & Prototipo web funcional con automatizaciones de motor. \\
			\midrule
			\textbf{Bloque 4: Fase de Validación} (Obj. Específico 4) & 13 & Nov 2026 & 13 & 4 semanas & Diseño de plan de pruebas. Ejecución de pruebas unitarias sobre el código PHP y funciones de motor. & Informes de cobertura de pruebas unitarias. \\
			\midrule
			\textbf{Bloque 4: Fase de Validación} (Obj. Específico 4) & 14 & Dic 2026 & 14 & 4 semanas & Comprobación de integración Frontend-Backend. Auditoría de vulnerabilidades (XSS, SQLi) en entorno simulado. & Reporte de integración y mitigación de fallas de seguridad. \\
			\midrule
			\textbf{Bloque 4: Fase de Validación} (Obj. Específico 4) & 15 & Ene 2027 & 15 & 4 semanas & Refinamiento del prototipo. Evaluación final de rendimiento con índices. Documentación de cierre del sistema. & Prototipo validado empíricamente, Manuales técnicos. \\
		\end{longtable}
	}
\end{landscape}

\section{Diagramas UML}



